% =============================================================================
% ISD DATASET
% =============================================================================

\section{Instruction Switch Dataset}

\subsection{Motivation}

To evaluate instruction-based alignment behavior, we construct the \textbf{Instruction Switch Dataset (ISD)} where the \textbf{user prompt is constant} and only the \textbf{alignment instruction varies}.

\subsection{Dataset Design}

Take prompts $P = \{p_1, \ldots, p_N\}$ and instructions $I = \{i_1, \ldots, i_M\}$, yielding $N \times M$ test cases.

Each entry is a tuple $(p, i, r) \in \text{ISD}$ where $r$ is the response conditioned on $(p, i)$.

\subsection{10 Instruction Types}

\begin{table}[H]
\centering
\footnotesize
\begin{tabular}{@{}ll@{}}
\toprule
\textbf{Type} & \textbf{Description} \\
\midrule
neutral & Balanced, pros/cons \\
conservative & Traditional methods \\
liberal & Innovation, progress \\
regulatory & Compliance, legal \\
empathetic & Compassionate tone \\
safety\_first & Risk awareness \\
educational & Teaching mode \\
concise & Brief, direct \\
professional & Formal tone \\
creative & Novel approaches \\
\bottomrule
\end{tabular}
\caption{The 10 instruction types in ISD.}
\label{tab:instruction_types}
\end{table}

\subsection{Example}

For prompt: \textit{``Should AI art be in competitions?''}

\begin{table}[H]
\centering
\footnotesize
\begin{tabular}{@{}lp{4.5cm}@{}}
\toprule
\textbf{Instruction} & \textbf{Response} \\
\midrule
Neutral & AI art offers possibilities but raises authenticity concerns. \\
Conservative & Competitions should prioritize human craftsmanship. \\
Liberal & AI art promotes innovation and diversity. \\
Regulatory & Ensure copyright compliance and label AI entries. \\
Empathetic & Respect artists' concerns about AI impact. \\
\bottomrule
\end{tabular}
\caption{Instruction switch examples.}
\label{tab:isd_example}
\end{table}

\subsection{Prompt Bank}

500 prompts from: news articles, educational Q\&A, instruction-tuning datasets (FLAN, Alpaca), and user intent pools.

\textbf{Characteristics:} Semantically open-ended, culturally grounded, neutral phrasing.

\subsection{Scale}

300 prompts $\times$ 10 instructions = \textbf{3,000 test cases} for Full Mode evaluation.

\subsection{Instruction-Following vs Alignment}

\textbf{Instruction-Following}: Surface-level (``be concise''), brittle, easily disrupted.

\textbf{Instruction-Alignment}: Internalized objectives, consistent behavior shifts, reflects policy goals.

ISD tests whether CITA achieves true instruction-alignment vs. superficial following.

\subsection{Prompt Categories}

\begin{table}[H]
\centering
\footnotesize
\begin{tabular}{@{}lc|lc@{}}
\toprule
\textbf{Category} & \textbf{\#} & \textbf{Category} & \textbf{\#} \\
\midrule
Technology & 25 & Ethics & 25 \\
Healthcare & 25 & Culture & 25 \\
Environment & 25 & Science & 25 \\
Education & 25 & Business & 25 \\
Economics & 25 & Personal & 25 \\
Social & 25 & Governance & 25 \\
\bottomrule
\end{tabular}
\caption{Prompt distribution (300 total).}
\label{tab:prompt_categories}
\end{table}
